% sections/07_gradient.tex:28-33
\begin{proposition}[Hidden-key computational hardness]\label{prop:hiddenkey}
Under the PRF assumption of Theorem~\ref{thm:prf}, no probabilistic polynomial-time learner independent of $s$, with a polynomial-time evaluable output and only polynomially many labeled samples from the uniform distribution on $\ell_0$, has average error at most $1/2-\gamma(n)$ for uniform $s$, where $\gamma(n)\ge1/\operatorname{poly}(n)$. More precisely, a learner with $q(n)$ samples and advantage $\gamma(n)$ gives a PRF distinguisher of advantage at least
\[
 \gamma(n)-\frac{q(n)}{2N}.
\]
\end{proposition}

% appendices/network_proofs.tex:13-15
\begin{proof}[Proof of \cref{prop:hiddenkey}]
Given an oracle function, generate the learner's independent uniform point samples on $\ell_0$ and label them by its oracle values at the corresponding incidence encodings. Run the learner without giving it a key. Draw an independent uniform test point on $\ell_0$, query its label, and output one if the predictor agrees. In the PRF case acceptance is at least $1/2+\gamma(n)$. In the random-function case, conditional on the training transcript, a previously unqueried point has an independent fair label. A fresh test point repeats a training point with probability at most $q(n)/N$. Acceptance is therefore at most $1/2+q(n)/(2N)$. The simulation and one output evaluation are polynomial-time. Since $N=2^{(n-1)/3}$, the collision term is negligible, contradicting PRF security for inverse-polynomial $\gamma$.
\end{proof}
